\paragraph{PageNavigation}
\textbf{Componente padre}:
\begin{itemize}
    \item DatasetContentPage;
    \item TestPage;
    \item TestComparisonPage;
    \item TestPageView;
    \item ComparisonPageView
\end{itemize}
\textbf{Elementi chiave}:
\begin{itemize}
    \item \textbf{PreviousButton}: pulsante per navigare alla pagina precedente;
    \item \textbf{NextButton}: pulsante per navigare alla pagina successiva;
    \item \textbf{PageInput}: campo per inserire direttamente il numero di pagina;
    \item \textbf{PageIndicator}: mostra la pagina corrente e il numero totale;
\end{itemize}
\textbf{Output}:
\begin{itemize}
    \item \textbf{next}: evento emesso quando si clicca il pulsante per la pagina successiva;
    \item \textbf{previous}: evento emesso quando si clicca il pulsante per la pagina precedente;
    \item \textbf{pageChange}: evento emesso quando si cambia pagina tramite il campo di input;
\end{itemize}
\paragraph{SearchBar}
\textbf{Componente padre}:
\begin{itemize}
    \item TestListPage;
\end{itemize}
\textbf{Elementi chiave}:
\begin{itemize}
    \item \textbf{InputField}: campo di input per inserire parole chiave da cercare;
    \item \textbf{SearchIcon}: icona cliccabile per avviare la ricerca;
\end{itemize}
\textbf{Input}:
\begin{itemize}
    \item \textbf{keywords}: parole chiave su cui effettuare la ricerca;
\end{itemize}
\textbf{Output}:
\begin{itemize}
    \item \textbf{search}: evento emesso quando viene inviata la conferma della ricerca;
    \item \textbf{clear}: evento emesso quando si clicca il pulsante di cancellazione;
\end{itemize} 
\paragraph{TestListPage}
\textbf{Modulo}:
\begin{itemize}
    \item TestModule, modulo funzionale dell'applicazione che raccoglie e organizza tutti i componenti, servizi e risorse legati all'esecuzione e gestione di test e alla visualizzazione dei risultati.    
\end{itemize}
\textbf{Elementi chiave}:
\begin{itemize}
    \item \textbf{SearchTest}: barra di ricerca per filtrare i test salvati, realizzata tramite il sottocomponente \textbf{SearchBar};
    \item \textbf{TestListContainer}: lista scorrevole dei test salvati, realizzata tramite il sottocomponente \textbf{TestListView};
    \item \textbf{UploadButton}: pulsante che permette di richiedere il caricamento di un test a partire da un file JSON, il caricamento viene poi effettuato tramite il sottocomponente \textbf{TestUploadForm}; 
\end{itemize}
\textbf{Servizio}
\begin{itemize}
    \item \textbf{TestService}: servizio iniettato tramite `@Injectable`, fornisce le operazioni per la gestione di un elemento della lista di test.
\end{itemize}
\paragraph{TestListView}
\textbf{Componente padre}:
\begin{itemize}
    \item TestListPage;
    \item TestPageView;
\end{itemize}
\textbf{Elementi chiave}:
\begin{itemize}
    \item \textbf{StatusMessage}: sezione testuale contenente il messaggio che viene visualizzato dall'utente nel caso in cui non ci siano test salvati nel sistema;
    \item \textbf{TestList}: lista scorrevole dei test salvati, permette la creazione di un sottocomponente \textbf{TestElement} per ogni test salvato (tramite la direttiva *\texttt{ngFor});
\end{itemize}
\textbf{Input}:
\begin{itemize}
    \item \textbf{tests}: lista di test salvati nel sistema;
\end{itemize}
\textbf{Output}:
\begin{itemize}
    \item \textbf{loadTest}: evento emesso quando riceve da \textbf{TestElement} la richiesta di caricare un test;
    \item \textbf{deleteTest}: evento emesso quando riceve da \textbf{TestElement} la richiesta di rimuovere un test;
    \item \textbf{renameTest}: evento emesso quando riceve da \textbf{TestElement} la richiesta di rinomina di un test;
\end{itemize}
\paragraph{TestElement}
\textbf{Componente padre}:
\begin{itemize}
    \item TestListView;
\end{itemize}
\textbf{Elementi chiave}:
\begin{itemize}
    \item \textbf{TestName}: sezione testuale contenente il nome del test;
    \item \textbf{TestDate}: sezione testuale contenente la data di esecuzione del test;
    \item \textbf{RenameButton}: pulsante per richiedere la rinominazione del test, questa viene poi effettuata tramite il sottocomponente \textbf{TestNameForm};
    \item \textbf{DeleteButton}: pulsante per richiedere l'eliminazione del test, questa viene poi confermata o annullata tramite il sottocomponente \textbf{ConfirmComponent};
    \item \textbf{LoadButton}: pulsante per richiedere il caricamento del test, se il test attualmente caricato presenta modifiche viene richiesta la conferma o l'annullamento della sovrascrittura tramite il sottocomponente \textbf{ConfirmComponent};
\end{itemize}
\textbf{Input}:
\begin{itemize}
    \item \textbf{dataset}: dataset associato al test;
    \item \textbf{llm}: LLM associato al dataset;
\end{itemize}
\textbf{Output}:
\begin{itemize}
    \item rename: evento emesso quando viene premuto il pulsante per richiedere la rinominazione del test;
    \item delete: evento emesso quando viene premuto il pulsante per richiedere l'eliminazione del test;
    \item load: evento emesso quando viene premuto il pulsante per richiedere il caricamento del test;
\end{itemize}
\paragraph{TestNameForm}
\textbf{Componente padre}:
\begin{itemize}
    \item TestPage;
    \item TestListPage;
\end{itemize}
\textbf{Elementi chiave}:
\begin{itemize}
    \item \textbf{NewNameLabel}: etichetta associata al campo di inserimento del nuovo nome;
    \item \textbf{NewNameInput}: campo di inserimento del nuovo nome da assegnare al test;
    \item \textbf{ConfirmButton}: pulsante per confermare il nuovo nome;
    \item \textbf{AbortButton}: pulsante per annullare la modifica;
\end{itemize}
\textbf{Input}:
\begin{itemize}
    \item \textbf{currentName}: gli viene passato il nome del test nel caso non sia temporaneo;
\end{itemize}
\textbf{Output}:
\begin{itemize}
    \item \textbf{confirmRename}: evento emesso quando viene premuto il pulsante per confermare l'assegnazione o modifica del nome del test;
    \item \textbf{cancelRename}: evento emesso quando viene premuto il pulsante per annullare l'assegnazione o modifica del nome del test;
\end{itemize}
\paragraph{TestPage}
\textbf{Modulo}:
\begin{itemize}
    \item TestModule, modulo funzionale dell'applicazione che raccoglie e organizza tutti i componenti, servizi e risorse legati all'esecuzione e gestione di test e alla visualizzazione dei risultati.    
\end{itemize}
\textbf{Elementi chiave}:
\begin{itemize}
    \item \textbf{NameText}: sezione testuale contenente il nome del test caricato;
    \item \textbf{TemporaryLabel}: sezione testuale contenente il messaggio che informa l'utente che il test caricato è temporaneo;
    \item \textbf{SummarySection}: sezione di visualizzazione dei dati sintetici del test, realizzata tramite il sottocomponente \textbf{TestIndex};
    \item \textbf{ComparisonButton}: pulsante che permette di richiedere il confronto tra il test corrente e un altro test salvato, la selezione del test da confrontare viene poi effettuata tramite il sottocomponente \textbf{TestSelectionComparison};
    \item \textbf{SaveButton}: pulsante che permette il salvataggio del test se è temporaneo, viene successivamente richiesto di assegnargli un nome tramite il sottocomponente \textbf{TestNameForm};
    \item \textbf{ResultSection}: sezione di visualizzazione dei risultati del test e del grafico di rappresentazione, realizzata tramite il sottocomponente \textbf{TestPageView};
    \item \textbf{ErrorMessage}: messaggio di errore che viene visualizzato in seguito a un errore nell'ottenimento o nel salvataggio dei dati, realizzato tramite il sottocomponente \textbf{MessageComponent};
\end{itemize}
\textbf{Servizio}
\begin{itemize}
    \item \textbf{TestService}: servizio iniettato tramite `@Injectable`, fornisce le funzionalità principali per la gestione di un test.
\end{itemize}
\paragraph{TestPageView}
\textbf{Componente padre}:
\begin{itemize}
    \item TestPage;
\end{itemize}
\textbf{Elementi chiave}:
\begin{itemize}
    \item \textbf{NavBar}: sezione che permette la navigazione tra le pagine del test, realizzata tramite il sottocomponente \textbf{PageNavigation};
    \item \textbf{Results}: pagina di visualizzazione dei risultati del test, realizzata tramite il sottocomponente \textbf{ResultListView};
    \item \textbf{Graph}: sezione di visualizzazione del diagramma di dispersione, realizzata tramite il sottocomponente \textbf{ScatterDiagram};
    \item \textbf{ErrorMessage}: messaggio di errore che viene visualizzato in seguito a un errore nell'ottenimento dei dati, realizzato tramite il sottocomponente \textbf{MessageComponent};
\end{itemize}
\textbf{Input}:
\begin{itemize}
    \item \textbf{Results}: array dei risultati da visualizzare;
\end{itemize}
\textbf{Output}:
\begin{itemize}
    \item \textbf{ChangePage}: evento emesso quando l'utente interagisce con `PageNavigation`;
\end{itemize}
\paragraph{TestIndex}
\textbf{Componente padre}:
\begin{itemize}
    \item TestPage;
    \item ComparisonIndex;
\end{itemize}
\textbf{Elementi chiave}:
\begin{itemize}
    \item \textbf{DataSummary}: sezione di visualizzazione dei dati sintetici del test, realizzata tramite i sottocomponenti: \textbf{CakeDiagram} e \textbf{BarDiagram};
    \item \textbf{ComparedSummary}: sezione di visualizzazione dei dati sintetici del test con cui si sta effettuando il confronto, realizzata tramite i sottocomponenti: \textbf{CakeDiagram} e \textbf{BarDiagram};
\end{itemize}
\paragraph{CakeDiagram}
\textbf{Componente padre}:
\begin{itemize}
    \item TestIndex;
\end{itemize}
\textbf{Elementi chiave}:
\begin{itemize}
    \item \textbf{PieChart}: sezione di visualizzazione del diagramma a torta;
    \item \textbf{ComparedPieChart}: sezione di visualizzazione del diagramma a torta del test con cui si sta effettuando il confronto;
    \item \textbf{PieElement}: singolo elemento del diagramma a torta che indica la percentuale di risposte corrette;
    \item \textbf{Legend}: legenda che descrive il significato di ogni colore usato nel diagramma;
\end{itemize}
\textbf{Input}:
\begin{itemize}
    \item \textbf{CorrectnessValue}: percentuale di risposte corrette;
\end{itemize}
\paragraph{BarDiagram}
\textbf{Componente padre}:
\begin{itemize}
    \item TestIndex;
\end{itemize}
\textbf{Elementi chiave}:
\begin{itemize}
    \item \textbf{BarChart}: sezione di visualizzazione del diagramma a barre;
    \item \textbf{ComparedBarChart}: sezione di visualizzazione del diagramma a barre del test con cui si sta effettuando il confronto;
    \item \textbf{BarElement}: singolo elemento del diagramma a barre che indica il numero di risposte per ogni intervallo (0-0.2, 0.2-0.4, 0.4-0.6, 0.6-0.8, 0.8-1);
    \item \textbf{BarLabel}: label che mostra l'intervallo che rappresenta l'elemento del diagramma;
    \item \textbf{Legend}: legenda che descrive il significato di ogni colore usato nel diagramma;
\end{itemize}
\textbf{Input}:
\begin{itemize}
    \item \textbf{distribution}: array con il numero di risultati per intervallo;
\end{itemize}
\paragraph{ScatterDiagram}
\textbf{Componente padre}:
\begin{itemize}
    \item TestPageView;
    \item ComparisonPageView;
\end{itemize}
\textbf{Elementi chiave}:
\begin{itemize}
    \item \textbf{ScatterChart}: sezione di visualizzazione del diagramma di dispersione;
    \item \textbf{ScatterElement}: singolo punto del diagramma di dispersione;
    \item \textbf{ComparedScatterElement}: singolo punto del diagramma di dispersione che rappresenta il test con cui si sta effettuando il confronto;
    \item \textbf{XLabel}: label che mostra il valore dell'asse X;
    \item \textbf{YLabel}: label che mostra il valore dell'asse Y;
    \item \textbf{Legend}: legenda che descrive il significato di ogni colore usato nel diagramma;
\end{itemize}
\textbf{Input}:
\begin{itemize}
    \item \textbf{ElementValues}: coordinate dei punti da rappresentare;
    \item \textbf{ComparedElementValues}: coordinate dei punti da rappresentare nel caso di confronto tra test;
\end{itemize}
\paragraph{ResultListView}
\textbf{Componente padre}:
\begin{itemize}
    \item TestPageView;
\end{itemize}
\textbf{Elementi chiave}:
\begin{itemize}
    \item \textbf{ResultList}: lista scorrevole dei risultati ottenuti dal test appartenenti a una pagina, permette la creazione di un sottocomponente \textbf{ResultElement} per ogni risultato (tramite la direttiva *\texttt{ngFor});
\end{itemize}
\textbf{Input}:
\begin{itemize}
    \item \textbf{results}: lista di risultati da visualizzare;
\end{itemize}
\paragraph{ResultElement}
\textbf{Componente padre}:
\begin{itemize}
    \item ResultListView;
\end{itemize}
\textbf{Elementi chiave}:
\begin{itemize}
    \item \textbf{Question}: sezione testuale contenente la domanda effettuata;
    \item \textbf{GivenAnswer}: sezione testuale contenente la risposta fornita dal LLM;
    \item \textbf{ExpectedAnswer}: sezione testuale contenente la risposta attesa;
    \item \textbf{Similarity}: sezione testuale contenente il grado di similarità tra la risposta fornita e quella attesa;
    \item \textbf{Correctness Label}: sezione testuale contenente il messaggio che indica se la risposta fornita è corretta o meno;
\end{itemize}
\textbf{Input}:
\begin{itemize}
    \item \textbf{question}: domanda effettuata;
    \item \textbf{givenAnswer}: risposta fornita dal LLM;
    \item \textbf{expectedAnswer}: risposta attesa;
    \item \textbf{similarity}: grado di similarità tra la risposta fornita e quella attesa;
    \item \textbf{correctness}: valore booleano che indica se la risposta fornita è corretta o meno;
\end{itemize}
\paragraph{TestSelectionComparison}
\textbf{Componente padre}:
\begin{itemize}
    \item TestPage;
\end{itemize}
\textbf{Elementi chiave}:
\begin{itemize}
    \item \textbf{SelectCompareList}: lista di elementi selezionabili usata per scegliere un test salvato con cui eseguire il confronto con il test attualmente visualizzato, la conferma o l'annullamento della selezione viene poi effettuata tramite il sottocomponente \textbf{ConfirmComponent};
\end{itemize}
    \textbf{Input}:
\begin{itemize}
    \item \textbf{dataset}: dataset associato al test;
    \item \textbf{llm}: LLM associato al dataset;
    \item \textbf{test}: test associato al dataset;
    \item \textbf{testList}: lista di test salvati nel sistema;
\end{itemize}
\textbf{Output}:
\begin{itemize}
    \item \textbf{compare}: evento emesso quando viene premuto il pulsante per richiedere il confronto tra test;
\end{itemize}
\paragraph{TestComparisonPage}  
\textbf{Modulo}:
\begin{itemize}
    \item TestModule, modulo funzionale dell'applicazione che raccoglie e organizza tutti i componenti, servizi e risorse legati all'esecuzione e gestione di test e alla visualizzazione dei risultati.    
\end{itemize}
\textbf{Elementi chiave}:
\begin{itemize}
    \item \textbf{Index}: sezione di visualizzazione dei dati sintetici del test, realizzata tramite il sottocomponente \textbf{ComparisonIndex};
    \item \textbf{CompareGraphs}: sezione di visualizzazione del diagramma di dispersione, realizzata tramite il sottocomponente \textbf{ComparisonPageView};
    \item \textbf{ErrorMessage}: messaggio di errore che viene visualizzato in seguito a un errore nell'ottenimento dei dati, realizzato tramite il sottocomponente \textbf{MessageComponent};
\end{itemize}
\textbf{Servizio}
\begin{itemize}
    \item \textbf{TestService}: servizio iniettato tramite `@Injectable`, espone i metodi per il confronto tra test.
\end{itemize}
\paragraph{ComparisonIndex} 
\textbf{Componente padre}:
\begin{itemize}
    \item TestComparisonPage;
\end{itemize} 
\textbf{Elementi chiave}:
\begin{itemize}
    \item \textbf{MetricComparison}: sezione di visualizzazione delle metriche di confronto tra i due test, realizzata tramite il sottocomponente \textbf{TestIndex};
\end{itemize}
\textbf{Input}:
\begin{itemize}
    \item \textbf{TestIndex1}: metriche del primo test;
    \item \textbf{TestIndex2}: metriche del secondo test;
\end{itemize}
\paragraph{ComparisonPageView}  
\textbf{Componente padre}:
\begin{itemize}
    \item TestComparisonPage;
\end{itemize}
\textbf{Elementi chiave}:
\begin{itemize}
    \item \textbf{ComparisonScatter}: sezione di visualizzazione del diagramma di dispersione, realizzata tramite il sottocomponente \textbf{ScatterDiagram};
    \item \textbf{ComparisonResults}: sezione di visualizzazione della lista dei risultati confrontati, realizzata tramite il sottocomponente \textbf{ComparisonListView};
    \item \textbf{NavBar}: sezione che permette la navigazione tra le pagine del test, realizzata tramite il sottocomponente \textbf{PageNavigation};
    \item \textbf{ErrorMessage}: messaggio di errore che viene visualizzato in seguito a un errore nell'ottenimento dei dati, realizzato tramite il sottocomponente \textbf{MessageComponent};
\end{itemize}
\textbf{Input}:
\begin{itemize}
    \item \textbf{comparisonData}: oggetto contenente le differenze per singola domanda;
\end{itemize}
\textbf{Output}:
\begin{itemize}
    \item \textbf{ChangePage}: evento emesso quando l'utente interagisce con `PageNavigation`;
\end{itemize}
\paragraph{ComparisonListView}
\textbf{Componente padre}:
\begin{itemize}
    \item ComparisonPageView;
\end{itemize}
\textbf{Elementi chiave}:
\begin{itemize}
    \item \textbf{ComparisonList}: lista scorrevole dei risultati confrontati, permette la creazione di un sottocomponente \textbf{ComparisonElement} per ogni risultato (tramite la direttiva *\texttt{ngFor});
\end{itemize}
\textbf{Input}:
\begin{itemize}
    \item \textbf{results1}: lista di risultati del primo test;
    \item \textbf{results2}: lista di risultati del secondo test;
\end{itemize}
\paragraph{ComparisonElement}
\textbf{Componente padre}:
\begin{itemize}
    \item ComparisonListView;
\end{itemize}
\textbf{Elementi chiave}:
\begin{itemize}
    \item \textbf{Question}: sezione testuale contenente la domanda effettuata;
    \item \textbf{GivenAnswer1}: sezione testuale contenente la risposta fornita dal LLM del primo test;
    \item \textbf{GivenAnswer2}: sezione testuale contenente la risposta fornita dal LLM del secondo test;
    \item \textbf{ExpectedAnswer1}: sezione testuale contenente la risposta attesa del primo test;
    \item \textbf{ExpectedAnswer2}: sezione testuale contenente la risposta attesa del secondo test;
    \item \textbf{Similarity1}: sezione testuale contenente il grado di similarità tra la risposta fornita e quella attesa del primo test;
    \item \textbf{Similarity2}: sezione testuale contenente il grado di similarità tra la risposta fornita e quella attesa del secondo test;
    \item \textbf{Correctness Label1}: sezione testuale contenente il messaggio che indica se la risposta fornita è corretta o meno del primo test;
    \item \textbf{Correctness Label2}: sezione testuale contenente il messaggio che indica se la risposta fornita è corretta o meno del secondo test;
\end{itemize}
\textbf{Input}:
\begin{itemize}
    \item \textbf{question}: domanda effettuata;
    \item \textbf{givenAnswer1}: risposta fornita dal LLM del primo test;
    \item \textbf{givenAnswer2}: risposta fornita dal LLM del secondo test;
    \item \textbf{expectedAnswer1}: risposta attesa del primo test;
    \item \textbf{expectedAnswer2}: risposta attesa del secondo test;
    \item \textbf{similarity1}: grado di similarità tra la risposta fornita e quella attesa del primo test;
    \item \textbf{similarity2}: grado di similarità tra la risposta fornita e quella attesa del secondo test;
    \item \textbf{correctness1}: valore booleano che indica se la risposta fornita è corretta o meno del primo test;
    \item \textbf{correctness2}: valore booleano che indica se la risposta fornita è corretta o meno del secondo test;
\end{itemize}
